\section{Requisiti Non Funzionali}
\begin{itemize}
    \item \textbf{RNF1 - Compatibilità:} Il sistema è tenuto a garantire piena compatibilità con i principali browser: Mozilla Firefox (versione 115 o superiore), Chromium/Chrome (versione 110 o superiore), Microsoft Edge (versione 110 o superiore) e Safari (versione 16 o superiore). L'esperienza utente dovrà risultare coerente su Windows, macOS e Linux, con corretta visualizzazione delle interfacce e piena operatività delle funzionalità indipendentemente da browser e sistema operativo. È inoltre richiesta un'interfaccia responsive, utilizzabile anche su tablet con schermo di almeno 10 pollici.

    \item \textbf{RNF2 - Performance:} Tempi di risposta rapidi e prestazioni efficienti sono requisiti imprescindibili. La dashboard principale dovrà caricarsi entro 10 secondi nel 95\% delle richieste, mentre le ricerche degli edifici dovranno completarsi entro 5 secondi. L'aggiornamento dei dati in tempo reale avverrà con una latenza massima di 120 secondi dalla rilevazione da parte dei sensori, tramite polling automatico a intervallo configurabile; invece per i creazione dei grafici storici, anche su archi temporali di mesi, il tempo massimo di generazione e visualizzazione è impostato a 20 secondi. Tali soglie vengono rispettate anche nei picchi di utilizzo ossia fino a 20 utenti connessi simultaneamente.

    \item \textbf{RNF3 - Scalabilità:} L'architettura deve essere pensata per la scalabilità, così da monitorare in modo efficiente almeno 100 edifici pubblici e un totale di 500 sensori installati. Le prestazioni definite in RNF2 dovranno mantenersi anche all'aumentare del numero di edifici monitorati e del volume di dati storici accumulati. Il design deve essere capace di supportare l'espansione futura con l'aggiunta di nuovi edifici e sensori senza cambiamenti strutturali sul sistema.

    \item \textbf{RNF4 - Affidabilità:} L'affidabilità deve essere elevata infatti la disponibilità minima avrà valori del 99.5\%, ossia in un anno sarebbero 43 ore e 50 minuti di downtime massimo, riducendo al minimo sia i tempi di inattività pianificati sia quelli imprevisti. Su richiesta dell'amministratore dovrà essere possibile generare backup completi del database in formato JSON, con esclusione dei dati sensibili (hash delle password); la retention dei dati di rilevazione va resa configurabile, con un minimo di 1 giorno. In caso di guasto ai sensori o perdita temporanea di connessione, il sistema continuerà a operare sugli ultimi dati, segnalando in modo chiaro l'anomalia (stato offline dei sensori). Il ripristino completo del servizio dopo un failure non dovrà superare le 4 ore.

    \item \textbf{RNF5 - Sicurezza:} Il livello di sicurezza di dati e accessi deve essere elevato. Tutte le comunicazioni client-server avverranno tramite HTTPS con certificati SSL/TLS validi, mentre le password saranno salvate con algoritmi di hashing robusti (bcrypt o Argon2) e salt individuale. Il sistema inoltre implementerà validazione rigorosa degli input e sanitizzazione dei dati. Il tutto in piena conformità al GDPR e alla normativa italiana sulla protezione dei dati personali.

    \item \textbf{RNF6 - Usabilità:} Il livello di praticità dell'applicazione deve essere elevato: un operatore con competenze informatiche di base dopo meno di 3 ore di formazione dovrà saper utilizzare le funzionalità principali, ossia la dashboard, la ricerca edifici e la gestione allarmi . L'interfaccia quindi sarà intuitiva e coerente in ogni sezione, con diagrammi di lavoro ben definiti, icone e pulsanti facilmente riconoscibili.

    \item \textbf{RNF7 - Conformità normativa:} Le normative GDPR sulla protezione dei dati personali dovranno essere pienamente rispettate. Le comunicazioni client-server saranno cifrate tramite TLS 1.2 o versione successiva, mentre i dati sensibili a riposo nel database (Data at Rest) saranno protetti con crittografia standard (es. AES-256).

    \item \textbf{RNF8 - Accuratezza dei dati:} Raccolta e visualizzazione dei dati dovranno rispettare standard elevati di accuratezza. Le temperature rilevate dai termometri avranno una precisione di ±0.5°C, salvo guasti o malfunzionamenti dei sensori; i consumi energetici registrati dovranno raggiungere una precisione del 98\% rispetto ai valori reali dei contatori. I dati meteorologici, ottenuti tramite API esterna, fungono da fonte di supporto per il calcolo delle temperature storiche di riferimento, e tale fonte va indicata chiaramente all'interno dell'applicazione.

    \item \textbf{RNF9 - Manutenibilità:} Manutenzione e aggiornamenti futuri dovranno risultare agevoli fin dalla progettazione. Il codice sorgente sarà ben documentato tramite commenti e conforme a standard di coding consolidati, con architettura modulare che consenta di aggiornare o sostituire singoli componenti senza ripercussioni sull'intero sistema. È richiesta una suite di test automatici con copertura di almeno il 70\% del codice, oltre a documentazione tecnica completa — diagrammi architetturali, documentazione delle API, manuale di deployment — costantemente aggiornata e accessibile al team di sviluppo.
    
\end{itemize}